\chapter{1977年江苏南通地区（试点）}

\section{解答题}

\begin{problem}
\begin{enumerate}
\item
当 \(a\) 为何值时，\(2a>a\)；当 \(a\) 为何值时，\(2a=a\)；当 \(a\) 为何值时，\(2a<a\).

\item
当 \(x\) 为何值时，\(|x|=x\).
\end{enumerate}
\end{problem}

\begin{solution}
\begin{enumerate}
\item 由 \(2a>a\) 两边同时减去 \(a\)，得 \(a>0\)，所以当 \(a>0\) 时，\(2a>a\). 由 \(2a=a\) 两边同时减去 \(a\)，得 \(a=0\)，所以当 \(a=0\) 时，\(2a=a\). 由 \(2a<a\) 两边同时减去 \(a\)，得 \(a<0\)，所以当 \(a<0\) 时，\(2a<a\).

\item 当 \(x\geqslant0\) 时，\(|x|=x\)；当 \(x<0\) 时，\(|x|=-x\)，而 \(-x>x\)，不可能有 \(|x|=x\). 因此 \(x\geqslant0\).
\end{enumerate}
\end{solution}

\begin{problem}
写出下列函数的自变量 \(x\) 的取值范围（在实数范围内）：

\begin{enumerate}
\item \(y=\dfrac{1}{2-x}\).
\item \(y=\sqrt{2x-1}\).
\end{enumerate}
\end{problem}

\begin{solution}
\begin{enumerate}
\item 分式有意义的条件是分母不为零，即 \(2-x\neq0\)，所以 \(x\neq2\). 因此定义域为所有满足 \(x\in\R\) 且 \(x\neq2\) 的 \(x\).

\item 根式在实数范围内有意义的条件是被开方数不小于零，即 \(2x-1\geqslant0\)，解得 \(x\geqslant\dfrac12\). 因此定义域为 \(\left[\dfrac12,+\infty\right)\).
\end{enumerate}
\end{solution}

\begin{problem}
计算：\(\left(\dfrac{2}{3}\right)^2+\left(-5.8\right)^0-\left(\dfrac{8}{27}\right)^{-\frac23}+\left(\dfrac12\right)^{-2}\).
\end{problem}

\begin{solution}
因为 \(-5.8\neq0\)，所以 \(\left(-5.8\right)^0=1\). 又因为 \(\dfrac{8}{27}=\left(\dfrac23\right)^3\)，所以 \(\left(\dfrac{8}{27}\right)^{-\frac23}=\left(\dfrac23\right)^{-2}=\dfrac94\)，且 \(\left(\dfrac12\right)^{-2}=4\). 因而
\[
\left(\dfrac{2}{3}\right)^2+\left(-5.8\right)^0-\left(\dfrac{8}{27}\right)^{-\frac23}+\left(\dfrac12\right)^{-2}=\dfrac49+1-\dfrac94+4=\dfrac{115}{36}=3\dfrac7{36}
\]
\end{solution}

\begin{problem}
\(k\) 为何值时，方程 \(x^2-2kx+4=0\) 有两个相等的实数根？并求出这方程的根.
\end{problem}

\begin{solution}
方程为关于 \(x\) 的一元二次方程，\(x^2-2kx+4=0\) 有两个相等的实数根的充要条件是判别式为零. 因此
\[
\Delta=(-2k)^2-4\cdot1\cdot4=4k^2-16=0
\]
解得 \(k^2=4\)，即 \(k=2\) 或 \(k=-2\).

当 \(k=2\) 时，原方程为 \(x^2-4x+4=0\)，即 \((x-2)^2=0\)，所以相等的实数根为 \(x_1=x_2=2\). 当 \(k=-2\) 时，原方程为 \(x^2+4x+4=0\)，即 \((x+2)^2=0\)，所以相等的实数根为 \(x_1=x_2=-2\).
\end{solution}

\begin{problem}
已知 \(\sin\alpha=\dfrac35\)（\(90^\circ<\alpha<180^\circ\)），求 \(\cos\alpha\)，\(\tan\alpha\)，\(\cot\alpha\) 的值.
\end{problem}

\begin{solution}
因为 \(90^\circ<\alpha<180^\circ\)，所以 \(\alpha\) 是第二象限角，\(\cos\alpha<0\). 由同角三角函数的平方关系，
\[
\cos^2\alpha=1-\sin^2\alpha=1-\left(\dfrac35\right)^2=\dfrac{16}{25}
\]
故 \(\cos\alpha=-\dfrac45\). 于是 \(\tan\alpha=\dfrac{\sin\alpha}{\cos\alpha}=\dfrac{3/5}{-4/5}=-\dfrac34\)，\(\cot\alpha=\dfrac{\cos\alpha}{\sin\alpha}=\dfrac{-4/5}{3/5}=-\dfrac43\).
\end{solution}

\begin{problem}
试写出抛物线 \(y=1+4x-2x^2\) 的顶点坐标，对称轴方程和函数 \(y\) 的极值.
\end{problem}

\begin{solution}
将函数配方，得
\[
y=1+4x-2x^2=-2\left(x^2-2x-\dfrac12\right)=-2\left[(x-1)^2-\dfrac32\right]=-2(x-1)^2+3
\]
因此抛物线的顶点坐标为 \((1,3)\)，对称轴方程为 \(x=1\). 对任意实数 \(x\)，都有 \((x-1)^2\geqslant0\)，所以 \(y=-2(x-1)^2+3\leqslant3\)，且当 \(x=1\) 时等号成立. 故函数 \(y\) 的极大值为 \(3\)，无极小值.
\end{solution}

\begin{problem}
解方程：\(\dfrac{\lg(2x)}{\lg(x-4)}=2\).
\end{problem}

\begin{solution}
原方程有意义必须满足 \(2x>0\)，\(x-4>0\)，并且分母 \(\lg(x-4)\neq0\). 这些条件合并为 \(x>4\) 且 \(x\neq5\).

在此定义域内，方程两边同乘 \(\lg(x-4)\)，得 \(\lg(2x)=2\lg(x-4)\). 根据对数的幂运算法则，
\[
\lg(2x)=\lg\left((x-4)^2\right)
\]
由于常用对数函数在正数范围内是单调函数，故 \(2x=(x-4)^2\)，即
\(x^2-10x+16=0\)，\((x-8)(x-2)=0\).
所以候选值为 \(x=8\) 或 \(x=2\). 其中 \(x=2\) 不满足 \(x>4\)，舍去；\(x=8\) 满足定义域，且 \(\lg16/\lg4=2\). 因此原方程的解为 \(x=8\).
\end{solution}

\begin{problem}
求 \(200\) 以内的能被 \(6\) 整除的所有自然数的和.
\end{problem}

\begin{solution}
不超过 \(200\) 的正的 \(6\) 的倍数依次为 \(6\)，\(12\)，\(\ldots\)，\(6\times33=198\)，因为 \(6\times34=204>200\). 共有 \(33\) 项，构成首项为 \(6\)、末项为 \(198\) 的等差数列，因此它们的和为
\[
S=\dfrac{33(6+198)}{2}=\dfrac{33\times204}{2}=3366
\]
若将 \(0\) 计入自然数，增加的 \(0\) 不改变和，故所求和仍为 \(3366\).
\end{solution}

\begin{problem}
已知正四棱锥底面边长为 \(12\mathrm{~cm}\)，高为 \(2\sqrt7\mathrm{~cm}\)，求它的侧棱.
\end{problem}

\begin{solution}
如图，\(P-ABCD\) 是正四棱锥，\(O\) 是正方形 \(ABCD\) 的中心，\(PO\) 是棱锥的高，所以 \(PO\perp\) 平面 \(ABCD\)，且 \(PO=2\sqrt7\mathrm{~cm}\).

\solutionfigure{\bitmapfigure[max width=0.42\linewidth]{img/1977/jiangsu_nantong/q9_solution_fig1.png}}

在正方形 \(ABCD\) 中，\(AB=BC=12\mathrm{~cm}\). 由勾股定理，
\[
AC=\sqrt{AB^2+BC^2}=\sqrt{12^2+12^2}=12\sqrt2\mathrm{~cm}
\]
因为 \(O\) 是正方形的中心，所以 \(AO=\dfrac{AC}{2}=6\sqrt2\mathrm{~cm}\). 又因为 \(PO\perp\) 平面 \(ABCD\)，所以 \(PO\perp AO\)，三角形 \(PAO\) 是直角三角形. 因而
\[
PA=\sqrt{PO^2+AO^2}=\sqrt{(2\sqrt7)^2+(6\sqrt2)^2}=\sqrt{28+72}=\sqrt{100}=10\mathrm{~cm}
\]
所以它的侧棱为 \(10\mathrm{~cm}\).
\end{solution}

\begin{problem}
两根钢绳悬吊货物时，两条钢绳与铅垂线的夹角分别为 \(\alpha\) 和 \(\beta\)，货物的重量为 \(P\)，试用表达式分别写出钢绳所承受的拉力 \(F_1\) 和 \(F_2\).
\end{problem}

\begin{solution}
如图，货物静止，故两根钢绳的拉力与重力 \(P\) 的合力为零. 将两根拉力分别沿水平方向和铅垂方向分解，水平方向的分力相互平衡，铅垂方向的分力平衡重力，所以
\[
\begin{cases}
F_1\sin\alpha=F_2\sin\beta,\\
F_1\cos\alpha+F_2\cos\beta=P.
\end{cases}
\]

\solutionfigure{\bitmapfigure[max width=0.34\linewidth]{img/1977/jiangsu_nantong/q10_solution_fig1.png}}

将第二个等式两边乘以 \(\sin\beta\)，并利用第一个等式 \(F_2\sin\beta=F_1\sin\alpha\)，得
\[
P\sin\beta=F_1\sin\beta\cos\alpha+F_2\sin\beta\cos\beta=F_1\left(\sin\beta\cos\alpha+\sin\alpha\cos\beta\right)=F_1\sin\left(\alpha+\beta\right)
\]
所以
\[
F_1=\dfrac{P\sin\beta}{\sin\left(\alpha+\beta\right)}
\]
同理，将第二个等式两边乘以 \(\sin\alpha\)，并利用 \(F_1\sin\alpha=F_2\sin\beta\)，可得
\[
P\sin\alpha=F_2\left(\sin\beta\cos\alpha+\sin\alpha\cos\beta\right)=F_2\sin\left(\alpha+\beta\right)
\]
因此
\[
F_2=\dfrac{P\sin\alpha}{\sin\left(\alpha+\beta\right)}
\]
\end{solution}
